\renewcommand{\chaptertitle}{The Tools}

%  Make first page footer:
\fancypagestyle{firststyle}{%
\fancyhf{}% Clear header/footer
\fancyhead{}
\fancyfoot[L]{{\footnotesize
              %% We toggle between these:
              % Manuscript submitted for review.\\
              \authorname{}~\authorpatp{}.  ``\chaptertitle{}''.  {\it Nockonomicon} (2025):  75–$x$.
              }}
}
%  Arrange subsequent pages:
\fancyhf{}
\fancyhead[LE]{{\gothicfont Nockonomicon}}
\fancyhead[RO]{\chaptertitle{}}
\fancyfoot[LE,RO]{\thepage}

\chapter{\chaptertitle}

\thispagestyle{firststyle}


A Nock noun is latent until \texttt{*} evaluates it.  Nock formulas are constructed, applied as discrete events to a subject, and result in new nouns.  These are then used as subjects or formulas for further evaluation.  From a pragmatic standpoint, we need a platform to evaluate subject--formula pairs (which we will call a runtime) and a set of design patterns to build new subjects and formulas for evaluation.

\section{A Runtime Evaluation Layer}

\section{Executable Nouns}

An attempt may be made to evaluate any noun as a Nock formula; obviously most formulas are malformed and will crash.  Executable nouns are also only valid with reference to a given subject:  \lstinline[style=inlinecode]{*[42 [0 2]]} will fail because the subject 42 is an atom, and so cannot be addressed with the axis 2.

The most common way to evaluate Nock is for the formula to access a noun in the subject, modify it in a characteristic way, and then evaluate the resulting modified noun.  We call this stored pattern a gate, the most common and convenient way to implement a function or a subroutine in Nock.  The gate is stored with an opcode \lstinline[style=inlinecode]{1} so it can be retrieved with opcode \lstinline[style=inlinecode]{9} and manipulated as a constant noun prior to execution.

For instance, this Nock program makes a gate that returns its argument incremented by \lstinline[style=inlinecode]{1}:

\begin{lstlisting}[style=listingblock]
[8 [1 0] [1 4 0 6]]
\end{lstlisting}

It composes a constant \lstinline[style=inlinecode]{0} (default input) with the pinned gate formula \lstinline[style=inlinecode]{[1 4 0 6]}, which retrieves the argument (axis 6) and applies the increment opcode \lstinline[style=inlinecode]{4} to the retrieved value.  The \lstinline[style=inlinecode]{[0 1]} tail is expanded by the cell distribution rule so that the resulting Nock formula \lstinline[style=inlinecode]{[4 0 6]} is actually evaluated.

One way to invoke it with an opcode \lstinline[style=inlinecode]{9} looks like this:

\begin{lstlisting}[style=listingblock]
:: compose
[7
   :: pin a gate with default sample argument 0
   [8 [1 0] [1 4 0 6] 0 1]
   :: pin access to the subject locally
   [8 [0 1]
      :: invoke the gate with the variable argument
      [9 2
         10 [6 7 [0 3] 1 5] 0 2]]
]
\end{lstlisting}

The meat of the pattern lies in the opcode \lstinline[style=inlinecode]{9} expression, which we may gloss as follows:

\begin{enumerate}
  \item  Retrieve the gate (using opcode \lstinline[style=inlinecode]{0}).
  \item  Modify the gate via opcode \lstinline[style=inlinecode]{10} by replacing the \lstinline[style=inlinecode]{0} pin with the argument, in this case \lstinline[style=inlinecode]{5}.
  \item  Invoke the modified gate.
\end{enumerate}

You can expand it using the rules for opcodes 7, 8, and 9 (using 0 as an unused placeholder subject):

\begin{lstlisting}[style=listingblock]
*[a 7 b c]          *[*[a b] c]
*[a 8 b c]          *[[*[a b] a] c]
*[a 9 b c]          *[*[a c] 2 [0 1] 0 b]
\end{lstlisting}

One way to think of Nock functions is as hygienic macros in the Lisp tradition. The gate pattern is a way to implement a macro that takes one argument and produces a new noun based on it.

\subsection{Key Design Patterns}

executable nouns
building gates
event loop

\section{The Event Loop}
